\documentclass[11pt]{article}
\usepackage[margin=1in]{geometry}
\usepackage{amsmath,amssymb}
\usepackage[hidelinks]{hyperref}
\emergencystretch=3em

\newcommand{\Fp}{\mathcal F_p}
\newcommand{\M}{\mathcal M}
\newcommand{\N}{\mathcal N}

\title{Literature Status for the Canonical Inclusion Question in arXiv:1811.01265}
\author{Run fa\_banach\_001, lane 7}
\date{June 16, 2026}

\begin{document}
\maketitle

\section*{Status}

This note records that Question~6.2 of
Albiac--Ansorena--Cuth--Doucha, arXiv:1811.01265, is answered positively by the
later paper Albiac--Ansorena, arXiv:2603.28193.

No new theorem is claimed here. The point of this packet is to clear an
open-problem signal from the run queue and identify the exact later theorem
which settles it.

\section*{Source Question}

In Section~6 of arXiv:1811.01265, after proving that canonical maps need not be
isometric, the authors ask whether the weaker isomorphic-embedding conclusion
always holds:
\begin{quote}
Let $0<p<1$ and $\N\subseteq \M$ be two $p$-metric (or metric) spaces in
inclusion. Is the canonical linear map of $\Fp(\N)$ into $\Fp(\M)$ an
isomorphic embedding?
\end{quote}
This is labelled Question~6.2 in the published version cited by the later
paper.

\section*{Later Answer}

The introduction of arXiv:2603.28193 explicitly cites this as AACD2020,
Question~6.2, states that it had remained open beyond the metric case, and says
that the paper gives a positive solution in full generality.

The theorem-level statement is arXiv:2603.28193, Theorem~3.13:
\[
  \text{for every }0<p\leq 1\text{ there is }C(p)<\infty
\]
such that, whenever $(\M,d)$ is a pointed $p$-metric space and
$0\in\N\subset \M$, the canonical linearization
\[
  L[\N,\M;p]\colon \Fp(\N)\longrightarrow \Fp(\M)
\]
is an isomorphic embedding and
\[
  \|L[\N,\M;p]^{-1}\|\leq C(p).
\]

This directly answers the source question in the affirmative. The source
question also mentions ordinary metric spaces; these are included under the
source convention for $p$-metric spaces when $0<p\leq 1$, since if $d$ is a
metric then $d^p$ is again a metric.

\section*{Scope}

This packet only records the literature answer to the canonical inclusion
question. It does not address the neighboring classification question in
arXiv:1811.01265 asking which $p$-metric spaces force the canonical map to be
isometric for every subspace.

\section*{References}

\begin{enumerate}
\item F. Albiac, J. L. Ansorena, M. Cuth, and M. Doucha,
\emph{Lipschitz free $p$-spaces for $0<p<1$}, arXiv:1811.01265.
\item F. Albiac and J. L. Ansorena,
\emph{Lipschitz extensions into $p$-Banach spaces, and canonical embeddings of
Lipschitz-free $p$-spaces for $0<p<1$}, arXiv:2603.28193.
\end{enumerate}

\end{document}
